\sect{श्रावण-शुक्ल-पुत्रदा-एकादशी-माहात्म्यम्}
\label{sec:padma-shravana-shukla-putrada}

\ganapatyadiDhyanam
\padmaPuranam

\dnsub{अथ सप्तपञ्चाशत्तमोऽध्यायः॥५७}

\uvacha{युधिष्ठिर उवाच}

\twolineshloka
{श्रावणस्य सिते पक्षे किं नामैकादशी भवेत्}
{कथयस्व प्रसादेन ममाग्रे मधुसूदन}% ॥१॥

\uvacha{श्रीकृष्ण उवाच}

\twolineshloka
{शृणुष्वावहितो राजन् कथां पापहरां पराम्}
{यस्याः श्रवणमात्रेण वाजपेयफलं भवेत्}% ॥२॥

\twolineshloka
{द्वापरस्य युगस्याऽऽदौ पुरा माहिष्मती पुरे}
{राजा महीजिदाख्यातो राज्यं पालयति स्वकम्}% ॥३॥

\twolineshloka
{पुत्रहीनस्य तस्यैव न तद् राज्यं सुखप्रदम्}
{अपुत्रस्य सुखं नास्ति इहलोके परत्र च}% ॥४॥

\twolineshloka
{चिन्तयानस्य तस्यैवं कालो बहुतरो गतः}
{न प्राप्तश्च सुतो राज्ञा सर्वसौख्यप्रदो नृणाम्}% ॥५॥

\twolineshloka
{दृष्ट्वाऽऽत्मानं प्रवयसं राजा चिन्तापरोऽभवत्}
{तदाऽऽगतः प्रजामध्ये इदं वचनमब्रवीत्}% ॥६॥

\twolineshloka
{इह जन्मनि भो लोका न मया पातकं कृतम्}
{अन्यायोपार्जितं वित्तं क्षिप्तं कोशे मया न हि}% ॥७॥

\threelineshloka
{ब्रह्मस्वं देवद्रविणं न गृहीतं मया क्वचित्}
{न्यासापहारो न कृतः परस्य बहुपापदः}
{पुत्रवत् पालितो लोको धर्मेण विजिता मही}% ॥८॥

\twolineshloka
{दुष्टेषु पातितो दण्डो बन्धुपुत्रोपमेष्वपि}
{शिष्टास्तु पूजिता नित्यं न द्वेष्याश्च मया जनाः}% ॥९॥

\twolineshloka
{इत्येवं ब्रुवतो मार्गं धर्मयुक्तं द्विजोत्तमाः}
{कस्मान्मम गृहे पुत्रो न जातस्तद् विमृश्यताम्}% ॥१०॥

\twolineshloka
{इति वाक्यं द्विजाः श्रुत्वा सप्रजाः सपुरोहिताः}
{मन्त्रयित्वा नृपहितं जग्मुस्ते गहनं वनम्}% ॥११॥

\twolineshloka
{इतस्ततश्च पश्यन्त आश्रमानृषिसेवितान्}
{नृपतेर्हितमिच्छन्तो ददृशुर्मुनिसत्तमम्}% ॥१२॥

\twolineshloka
{तप्यमानं तपो घोरं निरालम्बं निरामयम्}
{निराहारं जितात्मानं जितक्रोधं सनातनम्}% ॥१३॥

\twolineshloka
{लोमशं धर्मतत्त्वज्ञं सर्वशास्त्रविशारदम्}
{दीर्घायुषं महात्मानं सकेशं ब्रह्मसम्मितम्}% ॥१४॥

\twolineshloka
{कल्पेकल्पे गते तस्य एकं लोम विशीर्यते}
{अतो लोमशनामायं त्रिकालज्ञो महामुनिः}% ॥१५॥

\twolineshloka
{तं दृष्ट्वा हर्षिताः सर्वे आजग्मुस्तस्य सन्निधिम्}
{यथान्यायं यथार्हं ते नमश्चक्रुर्यथोदितम्}% ॥१६॥

\threelineshloka
{विनयावनताः सर्वे ऊचुस्ते च परस्परम्}
{अस्मद् भाग्यवशादेव प्राप्तोऽयं मुनिसत्तमः}
{तास्तथा स प्रजा वीक्ष्य उवाच ऋषिसत्तमः}% ॥१७॥

\uvacha{लोमश उवाच}

\twolineshloka
{किमर्थमिह सम्प्राप्तः कथयध्वं सकारणम्}
{दर्शनाद्धृष्टमनसः स्तुवन्तश्चैव मां किमु}% ॥१८॥

\twolineshloka
{असंशयं करिष्यामि भवतां यद्धितं भवेत्}
{परोपकृतये जन्म मादृशानां न संशयः}% ॥१९॥

\uvacha{जना ऊचुः}

\twolineshloka
{श्रूयतामभिधास्यामो वयमागमकारणम्}
{संशयच्छेदनार्थाय तव सान्निध्यमागताः}% ॥२०॥

\twolineshloka
{पद्मयोनेः परतरस्त्वत्तः श्रेष्ठो न विद्यते}
{अतः कार्यवशात् प्राप्ताः समीपं भवतो वयम्}% ॥२१॥

\twolineshloka
{महीजिन्नाम राजाऽसौ पुत्रहीनोऽस्ति साम्प्रतम्}
{वयं तस्य प्रजा ब्रह्मन् पुत्रवत्तेन पालिताः}% ॥२२॥

\twolineshloka
{तं पुत्ररहितं दृष्ट्वा तस्य दुःखेन दुःखिताः}
{तपः कर्तुमिहायाता मतिं कृत्वा तु नैष्ठिकीम्}% ॥२३॥

\twolineshloka
{तस्य भाग्येन दृष्टोऽसि ह्यस्माभिस्त्वं द्विजोत्तम}
{महतां दर्शनेनैव कार्यसिद्धिर्भवेन्नृणाम्}% ॥२४॥

\threelineshloka
{उपदेशं वद मुने राज्ञः पुत्रो यथा भवेत्}
{इति तेषां वचः श्रुत्वा मुहूर्तं ध्यानमास्थितः}
{प्रत्युवाच मुनिर्ज्ञात्वा तस्य जन्म पुरातनम्}% ॥२५॥

\uvacha{लोमश उवाच}

\twolineshloka
{पुरा जन्मनि वैश्योऽयं धनहीनो नृशंसकृत्}
{वाणिज्यकर्मनिरतो ग्रामाद् ग्रामान्तरं भ्रमन्}% ॥२६॥

\twolineshloka
{ज्येष्ठे मासि सिते पक्षे दशमी दिवसे तथा}
{मध्यगे द्युमणौ प्राप्ते ग्रामसीम्नि जलाशयम्}% ॥२७॥

\twolineshloka
{कूपिकां सजलां दृष्ट्वा जलपाने मनोदधे}
{सद्यस्ततः सवत्सा च धेनुस्तत्र समागता}% ॥२८॥

\twolineshloka
{तृष्णातुरा निदाघार्ता तस्यामम्बु पपौ तु सा}
{पिबन्तीं वारयित्वा तामसौ तोयं पपौ स्वयम्}% ॥२९॥

\twolineshloka
{कर्मणा तेन पापेन पुत्रहीनो नृपो भवेत्}
{कस्यापि जन्मनः पुण्यात् प्राप्तं राज्यमकण्टकम्}% ॥३०॥

\uvacha{लोका ऊचुः}

\threelineshloka
{पुण्यात् पापं क्षयं याति पुराणे श्रूयते मुने}
{पुण्योपदेशं कथय येन पापक्षयो भवेत्}
{यथा भवत्प्रसादेन पुत्रो भवति भूपतेः}% ॥३१॥

\uvacha{लोमश उवाच}

\twolineshloka
{श्रावणे शुक्लपक्षे तु पुत्रदा नाम विश्रुता}
{एकादशी वाञ्छितदा कुरुध्वं तद् व्रतं जनाः}% ॥३२॥

\twolineshloka
{इति श्रुत्वा नमस्कृत्य मुनिमेत्य पुरं व्रतम्}
{यथाविधि यथान्यायं कृतं तैर्जागरान्वितम्}% ॥३३॥

\twolineshloka
{तस्य पुण्यं सुविमलं दत्तं नृपतये जनैः}
{दत्ते पुण्येऽथ सा राज्ञी गर्भमाधत्त शोभनम्}% ॥३४॥

\twolineshloka
{प्राप्तो प्रसवकाले सा सुषुवे पुत्रमूर्जितम्}
{श्रावणस्य सिते पक्षे कर्कटस्थे दिवाकरे}% ॥३५॥

\twolineshloka
{द्वादश्यां वासुदेवाय पवित्रारोपणं स्मृतम्}
{हेम-रौप्य-ताम्र-क्षौमैः सूत्रैः कौशेयपद्मजैः}% ॥३६॥

\twolineshloka
{कुशैः काशैश्च कार्पासैर्ब्राह्मण्या कर्तितैः शुभैः}
{स्नात्वा त्रिगुणितं सूत्रं त्रिगुणीकृत्य शोधयेत्}% ॥३७॥

\twolineshloka
{गोदोहान्तरिते काले पूर्वेद्युरधिवासनम्}
{ब्राह्मणेभ्यो नमस्कृत्य गुरुपादौ प्रणम्य च}% ॥३८॥

\twolineshloka
{गीतमङ्गलनिर्घोषः कुर्याज्जागरणं ततः}
{ब्राह्मणाः क्षत्रिया वैश्या भिल्लाः शूद्रास्तथैव च}% ॥३९॥

\twolineshloka
{स्वधर्मावस्थिताः सर्वे भक्त्या कुर्युः पवित्रकम्}
{ततः पवित्रं गुरवे दद्याद्वै विधिपूर्वकम्}% ॥४०॥

\twolineshloka
{ब्राह्मणान् वैष्णवांश्चैव गन्धपुष्पादिनाऽर्चयेत्}
{अतो देवेति मन्त्रेण द्विजो विष्णौ निवेदयेत्}% ॥४१॥

\twolineshloka
{शूद्रस्तु मूलमन्त्रेण यथा विष्णौ तथा शिवे}
{वर्षेवर्षे प्रकर्तव्यं पवित्रारोपणं नरैः}% ॥४२॥

\twolineshloka
{भुक्तिं मुक्तिं च इच्छद्भिः संसारे शोकसागरे}
{न करोति विधानेन पवित्रारोपणं तु यः}% ॥४३॥

\threelineshloka
{तस्य सांवत्सरी पूजा निष्फला वैष्णवस्य तु}
{श्रुत्वा माहात्म्यमेतस्या नरः पापात् प्रमुच्यते}
{इह पुत्रसुखं प्राप्य परत्र स्वर्गतिं भवेत्}% ॥४४॥

॥इति श्रीपाद्मे महापुराणे पञ्चपञ्चाशत्साहस्र्यां संहितायामुत्तरखण्डे उमापतिनारदसंवादान्तर्गतकृष्णयुधिष्ठिरसंवादे श्रावण-शुक्ल-पुत्रदा-एकादशी-माहात्म्यं नाम सप्तपञ्चाशत्तमोऽध्यायः॥५७॥

\genMangalaShloka

\hyperref[sec:ekadashi_mahatmyam_padma_puranam]{\closesub}
\clearpage

